\documentclass[12 pt, letterpaper]{hmcpset}
\usepackage{hw-preamble}

\name{Jayden Thadani}
\class{MATH 168 --- Algorithms}
\assignment{Final exam review assignment}
\duedate{5th December 2025}

\begin{document}

\begin{problem}[Exercise 1: Course reflection]
The following several questions concern your experience in CS 140. Thoughtful responses written in complete sentences will receive full credit, whether or not they express positive views of the course. 

\emph{Note: this reflection exercise is not anonymous, though we appreciate your honesty. You will have the opportunity to provide anonymous feedback on the course and the Profs through the course evaluation form.}

\begin{enumerate}
    \item What are you most proud of accomplishing during this course?
    \item What was your favorite part of the course? (This can refer to specific lectures, problems, or modules or to broader categories of activities in the course, e.g. grutoring or office hours.)
    \item What was your \emph{least} favorite part of the course?
    \item Is there anything you wish we'd had more time to cover, and if so, what would you have liked to have seen?
    \item What's one piece of advice you'd like to share with future students in this course? (This could be something you are happy you did, or something you wish you had done.)
\end{enumerate}
\end{problem}

\begin{solution}
	\begin{enumerate}
		\item I'm most proud of getting better at solving combinatorial/graph-theoretic problems. I'm also proud of myself for accepting that I needed grutoring help.

		\item My favourite part of the course was the $\mathsf{NP}$-completeness section.

		\item I really didn't like having to write pseudocode for every use it or lose it/dynamic programming problem. It took painfully long and I didn't feel like I gained anything from having done it

		\item I wish we could have spent more time on fairness and fair division problems. I think some of the topological combinatorics involved is really interesting (Sperner's lemma for fair division, for example).

		\item I would tell students to go to grutoring! It's really helpful and really rewarding.
	\end{enumerate}
\end{solution}

\pagebreak

\begin{problem}[Exercise 2: Structured review]
The following attestations are designed to reward you for taking the time to thoroughly review the course content in advance of the final. For several attestations, we ask for cursory evidence, but for the most part they work on the honor system. Attesting ``I did review activity $X$'' without actually doing it is a violation of the HMC Honor Code. 

Because these are optional exercises, you're welcome to do one or two of the parts without doing all three.

\begin{enumerate}
    \item (10 points) Work your way through the Final Exam Topic List posted on Canvas as follows:
    \begin{itemize}
        \item Flag any items that seem unfamiliar, that you remember to be difficult, or that you'd like to review. 
        \item Do a second pass through the flagged items to identify which topics are highest priority and what materials you'd like to review (HW? Slides? Videos?) 
    \end{itemize}

    \item (10 points) Work your way through your graded HW and the corresponding solutions. (We'll attempt to publish HW9 grades before the exam. We won't have time to grade HW10 before the exam, but we'll publish solutions.) No need to show your work.
    \begin{itemize}
        \item Flag problems and subquestions that you struggled with, remember finding confusing, or solved but would like to review (for some of these questions, it's been a while!)
        \item Do a second pass through the flagged items to identify which issues can be resolved quickly and which need more attention. Identify which topics are highest-priority for review and what materials you'd like to review.
    \end{itemize}

    \item (10 points) After doing one or both of the previous parts, you now have a list of study topics and a rough priority ordering. Finally, make a study plan: decide what time you'd like to allocate for review, how you're going to review, if/when you'd like to attend grutoring review sessions, and when you plan to take the final.
    
    I attest that I've made a study plan as described above.
    \[
    	\textsc{Yes} \quad \quad \quad \quad \textsc{No}
    \]
\end{enumerate}
\end{problem}

\begin{solution}
	\begin{enumerate}
		\item I attest that I worked through the Topic List as described above.
			\[
				\textsc{Yes}
			\]
			Attach a photo, screenshot or other evidence of work (no need for this to be particularly tidy---all we want is some visual evidence that you did the above).
			\begin{figure}[H]
				\centering
				\includegraphics[width = 0.3 \textwidth]{figures/F25-Final-ReviewSheet-1.jpeg}
				\includegraphics[width = 0.3 \textwidth]{figures/F25-Final-ReviewSheet.jpeg}
				\caption{Two screenshots of my review of the topics list --- highlighted topics are things I particularly need to go over}
			\end{figure}

		\item I attest that I worked through past HW as described above.
			\[
				\textsc{Yes}
			\]
			Attach a photo, screenshot or other evidence of work (no need for this to be particularly tidy---all we want is some visual evidence that you did the above).

			I flagged the following problems in particular.
			\begin{itemize}
				\item All lined up --- high priority
				\item Touchy typesetting
				\item Sounds like a duck
				\item 3SAT is $\mathsf{NP}$-complete --- high priority
				\item Pe'er Haps proof of correctness
				\item Not throwing away $k$ shots --- high priority
			\end{itemize}

			\begin{figure}[H]
				\centering
				\includegraphics[width = 0.5 \textwidth]{figures/IMG_9104.jpeg}
				\caption{Me reviewing ``Not throwing away $k$-shots'' at my laptop}
			\end{figure}

		\item I attest that I've made a study plan as described above.
			\[
				\textsc{Yes}
			\]
	\end{enumerate}
\end{solution}

\pagebreak

\begin{problem}[Exercise 3: A reflective re-do]
This is an opportunity to spend time revisiting a HW problem. The intent is that you choose a problem that you really ``didn't get'' the first time through: a problem you skipped, struggled through, or attempted with an approach that just didn't work. (If \emph{no} problem falls in this category, congratulations---you should probably skip this question and relax!)

You're \textbf{invited to consult the HW solutions} as you go, but please don't just copy the answers here. This exercise is intended as an extended hints-type reflection: What went wrong? How would you fix it? What's your strategy for similar problems on the final?

For full credit, your answer should be self-contained enough that Prof. Tim can easily follow it using the HW solutions for reference (and without looking up your previous submission). Please be specific both about what went wrong and how you'd fix it.

\begin{enumerate}
    \item Which past HW problem are you choosing? If relevant, what parts or subparts? Briefly restate the problem in your own words.
    
    \item What went wrong the first time around? 
    
    (Unclear on the concepts involved? Unclear what the problem was asking? Understood the problem, but got stuck at a certain point? Ran out of time? Too much outside work or other pressure?)
    \item How would you fix it? 
    
    (If you misunderstood the problem, what do you know now? If you got stuck, does the posted HW solution make sense? If your solution needed fixing, where and what are those fixes, specifically?) 
    
    \item What would your strategy be if a similar problem came up on the final? 
    
    (Do you feel confident that you could handle it? If not, do you have a review plan and relevant material to work through to get prepared? If getting to the point where you feel fully prepared isn't realistic, what's your strategy for making the best of the situation?)
\end{enumerate}
\end{problem}

\begin{solution}
	\begin{enumerate}
		\item I'm redoing problem 3 on homework 4 --- ``Touchy typesetting''. The problem is to take an array of words and figure out where to add linebreaks so as to minimise a certain penalty function that depends on the number of unused lines.

		\item I wrote out something close to the correct solution for the problem several times but didn't realise it was correct --- I kept thinking I had come up with a counterexample. I had more than enough time to solve the problem, but I spent forever going back and forth over whether it was correct. I only realised what the correct solution was after the deadline had already passed.

		\item My solution was essentially to do use it or lose it on line breaking at each word. This was correct. If I could fix my solution, I'd try to recognise this earlier and then write it up properly.

		\item If a similar problem came up on the final, I would use the context of the problem to understand that it was use it or lose it. If I felt stuck on the specifics of the algorithm, I would sketch out the pseudocode and the dynamic programming details. Thus, once I figured out whether the algorithm was correct I wouldn't have to spend more time writing it up. Hopefully, this would allow me to complete the problem within the given time, or if not, get more partial credit.
	\end{enumerate}
\end{solution}

\end{document}
